\section*{Chapitre \ref{ch:b2:measurement-statistics} --- Statistique de la mesure chimique}

\begin{solution}{exo:b2:measurement-statistics:1}
$\bar V = 62.41/5 = \qty{12.482}{mL}$~; $\sum(V_i - \bar V)^2 = 0.00328$, $s = \sqrt{0.00328/4} \approx
\qty{0.029}{mL}$~; $s/\sqrt5 \approx \qty{0.013}{mL}$.
\end{solution}

\begin{solution}{exo:b2:measurement-statistics:2}
$t = 2.776$ pour 4 degrés de liberté~: $12.482 \pm 2.776 \times 0.0128$, soit $(12.48 \pm 0.04)$~mL.
\end{solution}

\begin{solution}{exo:b2:measurement-statistics:3}
Les incertitudes relatives de la masse et du volume valent 0.04~\% et 0.06~\%~; combinées
quadratiquement,
\[\frac{u(c)}{c} = \sqrt{\Bigl(\frac{0.0002}{0.5000}\Bigr)^2 + \Bigl(\frac{0.15}{250.0}\Bigr)^2}
\approx 0.07~\%.\]
\end{solution}

\begin{solution}{exo:b2:measurement-statistics:4}
$\bar x = 1.5$, $\bar y = 0.31$, $S_{xx} = 5$, $S_{xy} = 0.99$~: $b = 0.198$, $a = 0.31 - 0.198 \times 1.5
= 0.013$.
\end{solution}

\begin{solution}{exo:b2:measurement-statistics:5}
$|10.12 - 10.00|/(0.08/\sqrt5) = 0.12/0.036 \approx 3.4 > 2.776$~: l'écart est significatif, la
méthode est biaisée (d'environ \qty{+0.12}{mg/L}).
\end{solution}

\begin{solution}{exo:b2:measurement-statistics:6}
$\mathrm{pH} = -\log_{10}c$, $\mathrm d\,\mathrm{pH}/\mathrm dc = -1/(c\ln10)$~:
$u(\mathrm{pH}) = (u(c)/c)/\ln10 = 0.02/2.303 \approx 0.009$.
\end{solution}

\begin{solution}{exo:b2:measurement-statistics:7}
$x_{\mathrm{LD}} = 3 \times 0.0012/0.0098 \approx \qty{0.37}{\micro g/L}$~; $x_{\mathrm{LQ}} = 10
\times 0.0012/0.0098 \approx \qty{1.2}{\micro g/L}$.
\end{solution}

\begin{solution}{exo:b2:measurement-statistics:8}
La droite donne \qty{3.11}{mg/L} dans les solutions mesurées~; l'échantillon avait été dilué cinq fois~:
$3.11 \times 50.0/10.0 \approx \qty{15.6}{mg/L}$.
\end{solution}

\begin{solution}{exo:b2:measurement-statistics:9}
\qty{0.05}{mg/L} est la \omterm{def:b2:measurement-statistics:validation}{répétabilité}, \qty{0.15}{mg/L} la \omterm{def:b2:measurement-statistics:validation}{reproductibilité}. Chaque laboratoire a son
propre petit \omterm{def:b2:measurement-statistics:validation}{biais} (étalonnage, étalons, instrument)~; ces \omterm{def:b2:measurement-statistics:validation}{biais} diffèrent d'un laboratoire à l'autre et
s'ajoutent à la dispersion, si bien que la \omterm{def:b2:measurement-statistics:validation}{reproductibilité} est toujours la plus grande.
\end{solution}

\begin{solution}{exo:b2:measurement-statistics:10}
$\partial S/\partial a = 0$ s'écrit $\sum(y_i - a - bx_i) = 0$~: la somme des \omterm{def:b2:measurement-statistics:calibration}{résidus} est nulle. En
divisant par
$n$~: $\bar y = a + b\bar x$, donc $(\bar x, \bar y)$ est sur la droite.
\end{solution}

\begin{solution}{exo:b2:measurement-statistics:11}
En une étape~:
\[\sqrt{(0.006/1.000)^2 + (0.08/100.0)^2} \approx 0.61~\%.\]
Pour une étape au dixième~:
\[\sqrt{(0.02/10.00)^2 + (0.08/100.0)^2} \approx 0.22~\%,\]
et deux étapes combinées quadratiquement donnent $0.22\sqrt2 \approx 0.30~\%$. La voie en deux étapes est
deux fois plus précise~: la petite pipette domine l'incertitude de la voie en une étape.
\end{solution}

\begin{solution}{exo:b2:measurement-statistics:12}
$z = 1.6/\sqrt{0.4^2 + 0.5^2} \approx 2.5 > 2$~: les résultats ne sont pas compatibles~; l'un des
laboratoires (au moins) présente un \omterm{def:b2:measurement-statistics:validation}{biais}. Avec les incertitudes élargies ($k = 2$), le premier donne
$9.6 \pm 0.8$, qui contient 10~; le second
$11.2 \pm 1.0$, entièrement au-dessus de 10. Avant tout verdict sur l'eau, il faut lever le désaccord,
par exemple à l'aide d'un matériau de référence.
% ledger: who:lead
\end{solution}

\begin{solution}{pb:b2:measurement-statistics:1}
\textbf{1.} L'acide modifie l'atomisation et la réponse~; étalons et échantillons doivent avoir la même
matrice pour que la pente s'applique.
\textbf{2.} $\bar x = 10$, $\bar y = 0.1004$, $S_{xx} = 250$, $S_{xy} = 2.445$~: $b = 0.00978$ par
\unit{\micro g/L}, $a = 0.1004 - 0.0978 = 0.0026$.
\textbf{3.} $+0.0004$, $-0.0005$, $+0.0006$, $-0.0013$, $+0.0008$.
\textbf{4.} Les signes alternent, aucune courbure~: la droite convient sur 0--\qty{20}{\micro g/L}.
\textbf{5.} $s_{y/x} = \sqrt{3.1 \times 10^{-6}/3} \approx 0.00102$.
\textbf{6.} Le blanc (acide, eau, four) donne lui-même une petite absorbance.
\textbf{7.} La sensibilité~: l'absorbance par \unit{\micro g/L} de plomb.
\textbf{8.} $\bar y_0 = 0.1010$~; $s = 0.0030$.
\textbf{9.} $x_0 = (0.1010 - 0.0026)/0.00978 \approx \qty{10.06}{\micro g/L}$.
\textbf{10.} $s_{x_0} = (0.00102/0.00978)\sqrt{1/3 + 1/5 + (0.1010 - 0.1004)^2/(0.00978^2 \times 250)}
\approx 0.104 \times 0.730 \approx \qty{0.076}{\micro g/L}$.
\textbf{11.} $n - 2 = 3$~; $t = 3.182$.
\textbf{12.} $10.06 \pm 3.182 \times 0.076 = (10.06 \pm 0.24)$~\unit{\micro g/L}.
\textbf{13.} Le terme $1/m$ passe de 1 à $1/3$~: $s_{x_0}$ diminue d'environ un tiers.
\textbf{14.} $f = 50.50/50.0 = 1.010$~; $10.06 \times 1.010 \approx \qty{10.16}{\micro g/L}$.
\textbf{15.} $f = 1 + V_2/V_1$~: $u^2(f) = (u(V_2)/V_1)^2 + (V_2u(V_1)/V_1^2)^2 = (0.005/50.0)^2 +
(0.50 \times 0.03/2500)^2$, $u(f) \approx 1.0 \times 10^{-4}$, soit 0.01~\% en valeur relative.
\textbf{16.} Non~: 0.01~\% contre $0.076/10.06 \approx 0.75~\%$.
\textbf{17.} $s_{\mathrm{blanc}} \approx 0.00115$~; $x_{\mathrm{LD}} = 3 \times 0.00115/0.00978 \approx
\qty{0.35}{\micro g/L}$.
\textbf{18.} $10 \times 0.00115/0.00978 \approx \qty{1.2}{\micro g/L}$.
\textbf{19.} Oui, de loin.
\textbf{20.} Oui~: de 9.92 à \qty{10.40}{\micro g/L} dans l'eau du robinet.
\textbf{21.} $(0.104 - 0.0026)/0.00978 \times 1.010 \approx \qty{10.5}{\micro g/L}$, et, lue seule, un
«~au-dessus de la valeur guide~» affirmé avec assurance, que les données ne justifient pas.
\textbf{22.} La règle qui sert à construire l'intervalle capture la valeur vraie dans 95~\% des séries
auxquelles on l'applique~; ce n'est pas une probabilité portant sur cette valeur particulière.
\textbf{23.} Davantage de lectures répétées et davantage d'étalons vers \qty{10}{\micro g/L} (les termes
$1/m$ et
$1/n$), ou une méthode plus sensible~; l'intervalle ne se resserre que lentement, comme $1/\sqrt{\text{nombre}}$.
\textbf{24.} En analysant une eau de référence certifiée, ou en dopant l'échantillon par une quantité
connue de plomb et en vérifiant le taux de récupération.
\textbf{25.} Vers \qty{10}{\micro g/L}, l'incertitude des méthodes de routine est de quelques pour cent,
comme ici~: une valeur guide plus basse ne pourrait pas être contrôlée de façon fiable par des
laboratoires ordinaires.
\textbf{26.} $\boldsymbol{(10.16 \pm 0.24)}$~\unit{\micro g/L} à 95~\%~: l'intervalle contient
\qty{10}{\micro g/L}, si bien que la mesure ne montre ni que l'eau dépasse la valeur guide, ni qu'elle la
respecte avec une marge.
\end{solution}
